%% T3_localisation_paquet_ondes_FR.tex
\section{Localisation par paquets d’ondes (T3) : démonstration complète}

On suppose T2 établie et la formule explicite connue.

\subsection{Construction du paquet}
Soit $b\in C_c^\infty([-1,1])$, paire, $\int b=1$.
Pour $T\ge1$, posons :
\[
\widehat{\phi}_T(\gamma)
=\tfrac12\big(b(\sqrt{T}(\gamma-\gamma_0))+b(\sqrt{T}(\gamma+\gamma_0))\big).
\]
Alors $\widehat{\phi}_T$ est concentrée près de $\pm\gamma_0$ avec largeur $\sim T^{-1/2}$, et $\|\phi_T\|_2=1$.

\subsection{Terme archimédien \texorpdfstring{$\to0$}{to0}}
On a $K(t)=2\log|t|+O(|t|^{-2})$.
La convergence dominée (fonction majorante intégrable $h$) assure
\[
\frac{1}{2\pi}\int K(t)\widehat{\phi}_T(t)\,dt\to0.
\]

\subsection{Côté premiers \texorpdfstring{$\to0$}{to0}}
Par intégrations par parties, pour tout $M\ge1$ il existe $C_M$ tel que
\[
|\phi_T(x)|\le C_M\,T^{(M-1)/2}(1+|x|)^{-M}.
\]
Le nombre de $\xi_{p,k}\le A$ est fini, donc
\[
\sum_{p,k}\frac{\log p}{p^{k/2}}\,|\phi_T(\xi_{p,k})|\to0.
\]

\subsection{Contradiction dans le schéma}
En insérant dans la formule explicite,
\[
\sum_{\rho}\widehat{\phi}_T(\gamma_\rho)
=\frac{1}{2\pi}\int K\,\widehat{\phi}_T
+\sum_{p,k}\frac{\log p}{p^{k/2}}\phi_T(\xi_{p,k}).
\]
Le membre gauche $\to1$ (capture de masse), le droit $\to0$ : contradiction si un zéro $\rho_0$ est hors ligne.
